\documentclass[sigconf,nonacm]{acmart}

\usepackage{xcolor}

\newcommand{\todo}[1]{\textcolor{red}{\textbf{[TODO: #1]}}}
\newcommand{\citeTODO}[1]{\textcolor{orange}{\textbf{[CITE: #1]}}}
\newcommand{\question}[1]{\textcolor{purple}{\textbf{[QUESTION: #1]}}}
\newcommand{\result}[1]{\textcolor{blue}{\textbf{[RESULT: #1]}}}

\AtBeginDocument{%
  \providecommand\BibTeX{{\normalfont B\kern-0.5em{\scshape i\kern-0.25em b}\kern-0.8em\TeX}}}

\begin{document}
\raggedbottom

\title{Finding the Right Word? Evaluating Computational Lexical Suggestions for Creative Writing}

\author{Kathina Wardhani}
\affiliation{%
  \institution{The University of Sydney}
  \city{Sydney}
  \state{New South Wales}
  \country{Australia}
}

\begin{abstract}
\todo{Write this after the rest of the report is more mature}
\end{abstract}

\maketitle


\section{Introduction}

Finding the right word is a familiar part of writing. Writers use thesauruses not only to find exact synonyms, but also to search for words that are more specific, less clichéd, or carry different connotations \cite{gero2019stylistic}. Unexpected suggestions can also shape the direction of the writing itself \cite{gero2019stylistic}. Word search can therefore serve more than one purpose: it can help a writer refine an existing expression, but it can also expose alternatives that lead the writing somewhere new.

Computational approaches to word suggestion have largely developed around lexical substitution, where the task is to find alternatives for a target word that preserve its meaning and remain appropriate in context \cite{zhou2019bert,omarov2023grounding}. Even within this formulation, these requirements are not identical. A candidate that fits the surrounding sentence can diverge from the meaning of the original word, motivating methods that balance target semantics with contextual information \cite{zhou2019bert}. 

For creative writing, however, preserving the original expression is only part of the picture. Work on creativity support shows that generated material can contribute without functioning as a direct replacement. Singh et al. describe ``integrative leaps,'' where writers did additional work to incorporate generated material of varying semantic relevance into developing stories \cite{singh2022elephant}. The semantic distance and direction of word recommendations have also been shown to affect subsequent idea generation
\cite{wise2024sparking}. Together with observations of thesaurus use \cite{gero2019stylistic}, this suggests that a lexical suggestion can be valuable both for how well it fits what a writer has already expressed and for the alternatives it makes available to explore.

We therefore compare several ways of generating word suggestions for creative writing, from conventional lexical resources and a masked language model to instruction-tuned language models prompted for meaning-preserving or more exploratory alternatives. We consider both the properties of individual suggestions---how well they preserve the original meaning, how the substitution changes the likelihood of the sentence, how predictable the suggested word is in context, and how rare it is---and the diversity of the alternatives produced as a set. This allows us to compare methods without assuming that any one property defines a ``good'' suggestion.
% \question{Is ``linguistic plausibility'' the right term given $\Delta$PLL and surprisal? Does surprisal belongs here or under exploration?}

What a method suggests also depends on the context it receives. In lexical substitution, surrounding context is used to identify substitutes appropriate to a particular occurrence and to preserve sentence meaning \cite{zhou2019bert,seneviratne2022cilex}. Context-sensitive word search similarly uses surrounding language to rank alternatives for writers \cite{wiegmann2022language}. Creative writing also provides information beyond the immediate sentence, including the surrounding passage and examples of a writer's language. For the language-model methods, we test whether providing this broader document-level context changes the suggestions they produce.
\question{Replace the final sentence if context conditions test something narrower than sentence/passage/writer-language scope.}

Although these measures can characterize differences between methods, they cannot tell us how writers will respond to them. A semantically distant candidate might prompt a new direction or prove unhelpful, while a closely aligned candidate might provide the refinement a writer is looking for. Prior work also shows that writers can adapt and build upon generated material rather than accepting it outright \cite{singh2022elephant}. We therefore pair our computational comparison with a writer study to examine how suggestions with different profiles are used, rejected, adapted, and built upon.

Against this background, we ask:
\begin{itemize}
    \setlength{\itemsep}{2pt}
    \setlength{\parskip}{0pt}
    \setlength{\parsep}{0pt}
    \setlength{\topsep}{3pt}
    \setlength{\partopsep}{0pt}
    \item[\textbf{RQ1}] How do lexical suggestion methods differ in the
properties of individual suggestions and the sets of alternatives they produce?
    \item[\textbf{RQ2}] How does additional document-level context change the suggestions produced by language-model methods?
    \item[\textbf{RQ3}] How do writers use, reject, adapt, and build upon suggestions with different profiles?
\end{itemize}

% TODO: Add findings once the computational analysis and writer study are finalized.
\result{We find that (1) ...; (2) ...; and (3) ....}

% TODO: End with contributions following Wobbrock's Introduction structure.
\result{This work contributes (1) ...; (2) ...; and (3) ....}


\section{Related Work}

\subsection{Lexical Substitution}
% Traditional formulation
% Contextual lexical substitution
% BERT / neural / LLM approaches
% What counts as a "good" substitute

\subsection{Computational Support for Creative Writing}
% Creativity support tools
% Word-level recommendations
% Stolen Elephant / integrative leaps
% Phraselette
% Sparking Creativity

\subsection{Writer Agency, Intent, and Exploration}
% Suggestions do not have to be directly accepted to be useful
% Writer intent can evolve during interaction
% Agency / ownership / steering
% Connect this literature to alignment <-> exploration framing


\section{Computational Study}
\question{Formative study?}
\subsection{Lexical Suggestion Methods}
% Thesaurus baseline
% WordNet
% BERT MLM
% representative LLM conditions
% etc.

\subsection{Context Conditions}
% Sentence/local context
% broader document context
% author examples / other context conditions

\subsection{Evaluation Measures}
% Individual-suggestion properties:
%   semantic preservation: sentence STS + target-span similarity
%   change in sentence likelihood: delta PLL
%   contextual predictability: pseudo-surprisal / context gain
%   lexical rarity: word frequency
%
% Suggestion-set properties:
%   diversity
%
% Additional descriptive analyses:
%   uniqueness / overlap
%   yield / latency

\subsection{Experimental Setup}
% passages
% target-word selection
% candidate generation
% normalization/filtering
% repeated generation if applicable

\subsection{Results}
% RQ1: differences among methods
% RQ2: effects of context


\section{Writer Study}
\todo{Finalize study design with Katy.}
\subsection{Study Design}
% Controlled component
% Free/ecologically richer component?
% Within-subject / between-subject / mixed?
% Blinding / counterbalancing

\subsection{Participants}

\subsection{Materials and Interface}
% passages
% target words
% suggestion sets
% experimental interface

\subsection{Procedure}

\subsection{Measures and Analysis}
% accept / reject / adapt / build upon
% perceived usefulness
% agency/control
% qualitative rationales
% etc.

\subsection{Results}
% RQ3


\section{Discussion}

\subsection{Computational Trade-offs in Lexical Suggestion}

\subsection{The Role of Context}

\subsection{Implications for Creative-Writing Support}


\section{Limitations and Future Work}

\section{Conclusion}

\bibliographystyle{ACM-Reference-Format}
\bibliography{references}

\end{document}